% !TeX root = ../papers/nc-p3.tex
\section{平面局所環の定義式が四次以上のときの共通八点切断}
\label{sec:finite-high-order-common-octad}
\begin{theorem}\label{thm:finite-high-order-common-octad}
反対階数8で四辺が有限な候補のうち、既出の対称性または三文字の判定に入らない場合を考える。
平面局所点で $\dim\mathfrak m^j/\mathfrak m^{j+1}=j+1$ が $j=0,1,2,3$ に成り立てば、
その点は唯一の台である。各座標の接平面の式を $n_i$ とし、局所的に自明化する。
全ての $\operatorname{ann}(n_i^2)$ が同じイデアルとなり、これが定める長さ8の部分スキームは
各位置で三つの二次式の完全交叉である。この部分上で四つの座標は射影線型に同定される。
隣接する二つの接平面の積による辺の切断は長さ6、各座標の接平面切断は長さ7である。
\end{theorem}
\begin{proof}[二つの接平面と四つの関係]
平面局所環は二式完全交叉で、両式の位数が4以上なので局所交点数は16以上である。
他に未処理の局所因子があれば長さ7以上で全長20に反する。従って台は一点 $p$ である。
各側の共通四空間 $E$ の印での行列を $M$ とし、$M$ が第一対角成分だけを持つように選ぶ。
右下二行二列への射影を $C$ とする。印の二つの接座標を $x$、三座標の右・左核の一次部の
下二成分を $Rx,Lx$ とすると、どちらも二次元の同型である。
$A r=0$ および $l^t A=0$ から、$A$ の右下ブロックの一次部 $C_1(x)$ は
$C_1(x)Rx=(Lx)^tC_1(x)=0$（両式それぞれ零）を二次の多項式として満たす。
ここでは上・下の一次オフ対角成分と核の第一座標との積は三次以上で、二次の定義式がないことを使った。
非零の階数1以下の一次鉛筆なら共通像か共通余像を持つので、両核がともに動くことに反する。
従って $C_1=0$。$A$ が閉埋込みを与えるから、$M$ での階数一切断の接次元は少なくとも2である。
長さ4または正次元の切断は一点支持では既出の判定に入る。
従って両側とも接次元2の長さ3切断であり、定理\ref{thm:finite-marked-planar-triple}の $E_D$ 模型となる。
その $t$ 座標は非零の一次式で印の二次近傍の極大イデアルの二乗に入るので、その位置の接平面の式である。
行・列の先頭の射影座標も $M$ の核により対応する位置の接平面の式となる。

辺の二位置を $[s:x:n],[t:y:m]$ と書く。$x,y$ は各々二文字で、$n,m$ は接平面の式である。
両側の $E_D$ の行または列の評価から来る三関係ずつを取る。最初の組の下二式は
基底を選べば $xm=ny$ である。実際 $b$ の二文字、相手の下二文字、及び可逆な $D$ による式で、
接平面の式を加える二文字のリフトも選べる。他の組の下二式は同じ二空間を張る。
両組とも印の両方の評価で消える二空間に入り、その次元は正確に2だからである。
後者は辺の六式の一次部分が接次元2より階数4であることから従う。
それぞれの上式では相手側の下二文字と $b$ の二文字が非退化に対合する。
異なるリフトからの $xm,ny,nm$ の項を下二式と $s,t$ への一次式の加算で消せば、
四つの共有部分からの関係は
\[
                 xm=ny,\qquad sm=q(x,y),\qquad nt=r(x,y)
\]
となる。$q,r$ は二行二列の可逆双一次形式である。
二つは独立である。実際 $n=m=0$ の平面の積に制限した六式は、有限交叉のため少なくとも四つの独立式を必要とする。
従ってこの積で $q,r$ と残り二式は完全交叉で、長さは $\binom42=6$ である。

\end{proof}
\begin{proof}[八点の連結と二乗イデアル]
印のアフィン座標 $s=t=1$ で $m=q,n=r$ を消去する。辺の各射影は閉埋込みなので
残り二式の $y$ での微分は可逆である。形式的に $y=Tx+O(x^2)$ を解くと $T$ も可逆となる。
$xq-yr$ の三次初期式が恒等零であることから
\[
                     q(x,Tx)x=r(x,Tx)Tx .
\]
どちらの二次式も零なら、独立な $q,r$ が可逆な $T$ を通してともに唯一の交代形式の倍となり矛盾する。
よって一般の $x$ が $T$ の固有ベクトルであり $T=cI$, $c\ne0$。
従って $q=c r+\kappa[x,y]$、$\kappa\ne0$ である。
射影線型同型 $[s:x:n]\mapsto[s:cx:cn]$ のグラフに全六式を制限すると、
上の四式は $sn=c r(x,x)$ の一式となる。残り二式とともに有限な辺の部分を切るので
三二次の完全交叉、長さ8である。その支持も $p$ だけである。

この部分の局所イデアルを同定する。$y=y(x)$ を解いた二変数冪級数環で $d=y-cx$ と置く。
辺の二式は $Md$ で、$M=\kappa x[x,-]-r I$、$\det M=r^2$ である。
線型同型のグラフを課した部分のイデアルは $(d)$ である。
実際 $d=0$ とすれば $m=c n$ も従う。
二つの正則列 $(Md)\subset(d)$ の商の間で、$\det M$ 倍は後者の商を前者の商に単射し、
その像は $(d)$ の零化イデアルである。これはKoszul分解の比較で最上段が $\det M$ 倍で、
二次の局所双対が商の全射を単射に移すためである。
従って八点部分の定義イデアルは $\operatorname{ann}(r^2)=\operatorname{ann}(n^2)$ となる。
また $g=[x,y]$ は $xm=ny$ より $n,m$ の両方に消され、$m-cn=\kappa g$。
従って $m^2=c^2n^2$ がこの自明化で成り立つ。自明化を変えても二乗の主イデアルは変わらないから、
全辺で同じ八点部分を得る。四つの座標はこの部分の上で各辺の線型同型によって対応する。

最後に $n=0$ の切断を同じ二変数環で読むとそのイデアルは $(r,x_1g,x_2g)$ である。
$(r,g)$ の商は $n=m=0$ の長さ6切断であり、$g$ は一次元環 $k[[x_1,x_2]]/(r)$ の非零因子である。
従って $g$ の像を極大イデアル倍まで残した商の長さはちょうど一つ増えて7である。
\end{proof}
